% Article: Modelo Revista InFor
\documentclass{../../common-shared/cls/Article_Class_RevInfor}

% Search path for images
\graphicspath{{../../common-shared/img/}}

% Título do artigo (o nome próprio "Plan Ceibal" fica em itálico dentro do título,
% como no PDF oficial). A classe converte o título para CAIXA ALTA automaticamente.
\titulo{Contornos do currículo e a integração de tecnologias no contexto da política pública \textit{Plan Ceibal} desenvolvida no Uruguai}

% Dados do autor em partes. Deles a classe monta o nome por extenso do bloco de
% autoria (\nomeCompletoAutorRevInfor) e o autor-data do cabeçalho corrido das
% páginas ímpares ("SILVA, 2025").
\nomeAutorRevInfor{Jayson Magno da}
\sobrenomeAutorRevInfor{Silva}
\anoRevInfor{2025}   % ano da publicação -- ajustar no fechamento da edição

% Bloco de autoria no padrão Revista InFor:
%   \autorRevInfor{nome}{afiliação}{e-mail}{ORCID iD}
\autorRevInfor
    {\nomeCompletoAutorRevInfor}
    {Doutor em Educação. Professor Formador I, Pesquisador Colaborador na Universidade Federal do ABC (UFABC), São Paulo, Brasil.}
    {jayson.magno@ufabc.edu.br}
    {0000-0001-5858-2507}

% Banner de cabeçalho da Revista InFor, colado ao topo e às laterais do papel
% (full-bleed, margem zero). Impresso automaticamente acima do título.
\cabecalhoRevInfor{cabecalho-img-revistainfor.png}

\begin{document}

% ---------------------------------------------------------------------------
%  Elementos PRÉ-TEXTUAIS: resumo, palavras-chave, título e resumo em inglês.
%  Vêm depois do \begin{document} (o \maketitle é automático, disparado pela
%  classe) e ANTES do \textual, que abre o corpo do artigo.
%
%  O texto abaixo é o do artigo publicado em temp/10-23-PB.pdf, p. 8, para que
%  o PDF gerado possa ser comparado lado a lado com o oficial da revista.
% ---------------------------------------------------------------------------
\begin{resumo}
O presente artigo tem por objetivo refletir sobre a integração currículo e tecnologias digitais por meio de um retrato da pesquisa concluída que tomou como objeto a política pública \textit{Plan Ceibal} desenvolvida no contexto da República Oriental do Uruguai, sob o argumento do uso dos computadores portáteis na educação aliados ao potencial de mobilidade e conexão à Internet. O método está fundamentado na \textit{grounded theory} e teve como apoio o software \textit{NVivo 11}. Ao final é apresentado o \textit{constructo teórico} que decorre na emergência e análise dos achados, e conclui que a pesquisa apontou o papel preponderante e fundamental dos sujeitos na (re)construção de novos contextos experienciados com uso dos computadores portáteis aliados ao potencial de mobilidade e conexão à Internet, evidenciando indícios da integração entre currículo e tecnologias digitais e a emergência de \textit{web currículo}s.

% As palavras-chave saem em PARÁGRAFO PRÓPRIO, com o rótulo em negrito — o
% \palavraschaveRevInfor já fecha o parágrafo anterior (\par) antes de imprimir.
\palavraschaveRevInfor{Currículo e Tecnologias. Plan Ceibal. Computadores Portáteis. Tecnologias Digitais de Informação e Comunicação. Políticas Públicas.}
\end{resumo}

% Título em inglês, entre as palavras-chave e o cabeçalho "ABSTRACT".
\tituloestrangeiroRevInfor{Curriculum outlines on the context of the public policy Plan Ceibal developed in Uruguay}

\begin{resumo}[Abstract]
\begin{otherlanguage*}{english}
This article aims to reflect about curriculum integration and digital technologies through portrait of the completed research that took as object the public policy Plan Ceibal developed on context of the Oriental Republic of Uruguay, under the argument of the use of laptop in education allied to the potencial for mobility and Internet connection. The method is based on grounded theory and was supported by the NVivo 11 software. At the end it has the theoretical constructo that arises in the emergence and analysis of the findings and concludes that the research pointed out the preponderant and fundamental role of the subjects in the (re)constrution of new contexts experienced with the use of laptops combined with the potencial for mobility and Internet connection, showing evidence of integration between curriculum and digital technologies and the emergence of web curriculum.

\palavraschaveRevInfor[Keywords]{Curriculum and Technologies. Plan Ceibal. Laptops. Digital Information and Communication Technologies. Public Policy.}
\end{otherlanguage*}
\end{resumo}

\textual

\section{Introdução}
Este é o modelo refatorado para a Revista InFor. O conteúdo foi organizado para permitir múltiplos artigos no mesmo repositório compartilhando a mesma classe e imagens de cabeçalho.

Este é o modelo refatorado para a Revista InFor. O conteúdo foi organizado para permitir múltiplos artigos no mesmo repositório compartilhando a mesma classe e imagens de cabeçalho.


Este é o modelo refatorado para a Revista InFor. O conteúdo foi organizado para permitir múltiplos artigos no mesmo repositório compartilhando a mesma classe e imagens de cabeçalho.

Este é o modelo refatorado para a Revista InFor. O conteúdo foi organizado para permitir múltiplos artigos no mesmo repositório compartilhando a mesma classe e imagens de cabeçalho.

Este é o modelo refatorado para a Revista InFor. O conteúdo foi organizado para permitir múltiplos artigos no mesmo repositório compartilhando a mesma classe e imagens de cabeçalho.

\section{Desenvolvimento}
\lipsum[1-2]

\lipsum[3-4]

\lipsum[5-6]
\section{Considerações Finais}
\lipsum[7]

\lipsum[8]

\lipsum[9]

\postextual
% Shared bibliography (common to all articles) + this article's own references
\bibliography{../../common-shared/bib/references,references}

\end{document}
